% !TeX root = main.tex
\lecture{13}{Thu 5 Mar 2026 15:00}{Orbits I: Formalism of Orbits}

\section{Introduction to Orbits}
\begin{tcolorbox}
    \textbf{What is an Orbit?}

    \textbf{Orbit} $\equiv$ motion of an object in a plane (2D only) under the influence of a force. Generally, this is a closed repeating path (especially for this course) but this is not always true.
\end{tcolorbox}

For orbits, we limit ourselves to forces and potentials which are ``Central Conservative''.
\begin{itemize}
    \item Central Forces:
    \begin{itemize}
        \item A central force has the form $\underline{F} = f(r) \hat{\underline{r}}$. These forces act only in a radial direction and vary only as a function of radial distance. These forces conserve angular momentum.
        \item For a central force:
        \[
            \underline{F} = f(r) \hat{r} \implies \underline{\tau} = \underline{r} \times \underline{F} = f(r) \underline{r} \times \hat{r} = 0
        \]
        Hence there is no torque. Since torque is the derivative of angular momentum, angular momentum is conserved.
        \item Also, since $\underline{L} = \underline{r} \times \underline{p} = r \times m \dot{\underline{r}}$, $\underline{r} \times \dot{\underline{r}}$ is conserved. As $\underline{r}$ and $\ddot{\underline{r}}$ are both perpendicular to $\underline{L}$, they sit in the same 2D plane.
        
        This confirms that the orbits are exclusively 2D planes.
        
    \end{itemize}
    \item Conservative Forces:
    \begin{itemize}
        \item A conservative force is one that conserves total energy ($E_k + E_\text{potential}$). The net work done by returning to the start point via a closed path (most orbits) is zero.
        \[
            \oint \underline{F} \cdot d \underline{r} = 0
            \]
            \item For every conservative force, there exists a scalar potential which satisfies:
            \[
                \underline{F} = - \nabla V = - \frac{\partial V}{\partial r}
            \]
            We will therefore choose to work with forces and potentials interchangeably. Dealing with potentials is often easier due to them being a scalar.
    \end{itemize}
                
\end{itemize}

This includes all non-dissipative forces, i.e. gravity, electrostatics, the strong nuclear force etc. It excludes dissipative forces where energy is lost, such as friction.

\subsection{Polar Coordinates}
As we've established that coordinates exist purely in 2D, it makes the most sense to express them in polar coordinates, with the origin at the orbit centre.

Instead of being interested in unit vectors $\hat{x}, \hat{y}$ we're interested in unit vectors $\hat{\underline{r}}$ pointing radially outwards and $\hat{\underline{\theta}}$ pointing tangentially in the positive direction of motion (anticlockwise).

Given an arbitrary point:
\[
    (x, y) = (r, \theta)
\]

We can define:
\[
    \underline{r} = r \hat{\underline{r}} \tag{1}
\]

And:
\[
    \boxed{\hat{\underline{r}} = \cos \theta \hat{\underline{x}} + \sin \theta \hat{\underline{y}}}
\]
\[
    \boxed{\hat{\underline{\theta}} = - \sin \theta \hat{\underline{x}} + \cos \theta \hat{\underline{y}}}
\]

We want to be able to take their time derivatives:
\[
    \frac{d \hat{\underline{r}}}{dt} = \frac{d \theta}{dt} \frac{d}{d \theta} \left(\cos \theta \hat{\underline{x}} + \sin \theta \hat{\underline{y}}\right)
\]
\[
    = \dot{\theta} \left(- \sin \theta \hat{\underline{i}} + \cos \theta \hat{\underline{y}}\right)
\]
\[
    \boxed{\implies \frac{d \hat{\underline{r}}}{dt} = \dot{\theta} \hat{\underline{\theta}}} \tag{2}
\]

By the same logic:
\[
    \frac{d \hat{\underline{\theta}}}{dt} = \frac{d \theta}{dt} \frac{d}{d \theta} \left(- \sin \theta \hat{\underline{x}} + \cos \theta \hat{\underline{j}}\right)
\]
\[
    \implies \boxed{\frac{d \hat{\underline{\theta}}}{dt} = - \dot{\theta} \hat{\underline{r}}} \tag{3}
\]

These let us handle cases where the distances from the axes are changing (i.e. object in an elliptical orbit). The derivation is non-examinable and the results are in the formula sheet.

\section{Motion of an Object in Central Conservative Potential}
Assume an elliptical orbit. The origin of the coordinate system is at the left focus on the ellipse (the sun) and a planet is in orbit.

The planet sits at varying distance $r$ from the sun and has angle from the positive horizontal $\theta$.

\subsection{Planetary Velocity}
By definition:
\[
    \underline{v} = \dot{\underline{r}} = \frac{d}{dt} \underline{r}= \frac{d}{dt} \left(r \hat{\underline{r}}\right)
\]
\[
    \underline{v} = r \frac{d \hat{\underline{r}}}{dt} + \frac{dr}{dt \hat{\underline{r}}} = r \dot{\theta} \hat{\underline{\theta}} + \dot{r} \hat{\underline{r}}
\]
These two components give the tangential and then the radial velocity.

\subsection{Planetary Acceleration}
To find the acceleration, we take the time derivative again. Because $r, \dot{r}, \dot{\theta}$, and the unit vectors themselves are all functions of time, we use the chain rule with respect to $t$.
\[
    \underline{a} = \frac{d}{dt} (\dot{r} \hat{\underline{r}}) + \frac{d}{dt} (r \dot{\theta} \hat{\underline{\theta}})
\]
\[
    \underline{a} = (\ddot{r} \hat{\underline{r}} + \dot{r} \dot{\hat{\underline{r}}}) + (\dot{r} \dot{\theta} \hat{\underline{\theta}} + r \ddot{\theta} \hat{\underline{\theta}} + r \dot{\theta} \dot{\hat{\underline{\theta}}})
\]
Substituting the unit vector time derivatives, $\dot{\hat{\underline{r}}} = \dot{\theta} \hat{\underline{\theta}}$ and $\dot{\hat{\underline{\theta}}} = -\dot{\theta} \hat{\underline{r}}$:
\[
    \underline{a} = \ddot{r} \hat{\underline{r}} + \dot{r} (\dot{\theta} \hat{\underline{\theta}}) + \dot{r} \dot{\theta} \hat{\underline{\theta}} + r \ddot{\theta} \hat{\underline{\theta}} + r \dot{\theta} (-\dot{\theta} \hat{\underline{r}})
\]
And factorising by unit vectors:
\[
    \underline{a} = (\ddot{r} - r \dot{\theta}^2) \hat{\underline{r}} + (r \ddot{\theta} + 2 \dot{r} \dot{\theta}) \hat{\underline{\theta}}
\]
This gives us a radial component of acceleration and a tangential component of acceleration.

\section{Forces}
We know that $\underline{F} = m \underline{a}$ must apply, and we can resolve this out into radial and tangential components.

In the radial direction:
\[
    f(r) = m \left(\ddot{r} - r \dot{\theta}^2\right) \tag{4}
\]

As the force is central conservative, there can be no component in the $\hat{\underline{\theta}}$ direction. Hence:
\[
    r \ddot{\theta} + 2 \dot{r} \theta = 0 \tag{5}
\]

\section{Angular Momentum}
We again know that $\underline{L} = \underline{r} \times \underline{F} = m \underline{r} \times \underline{v}$. Let the modulus of the angular momentum vector being $l = |\underline{L}|$.

Only the tangential component of the velocity has a cross product with the radius, hence we have:
\[
    l = mr v_t
\]

Where $v_t$ is the tangential portion of the velocity. Inserting the tangential velocity:
\[
    l = mr \times r \dot{\theta} = mr^2 \dot{\theta}
\]

As an aside, noting that $dl/dt = mr^2 \ddot{\theta} + 2m r \dot{r} \dot{\theta}$:
\[
    \frac{dl}{dt} = mr \left(r \ddot{\theta} + 2 \dot{r} \dot{\theta}\right)
\]

This is equal to $mr F_T$. We have already shown that $F_T = 0$, hence $l$ is not time dependant, so is conserved.

\section{Total Energy}
As the field is conservative, it must be a constant equal to the sum of the kinetic and potential components.

\[
    E = E_k + E_p = \frac{1}{2}m \underline{v}^2 + V(r)
\]

Where $V(r)$ is the central potential where $\underline{F} = - \frac{\partial V}{\partial r} \underline{\hat{r}}$. Using the known value for velocity:
\[
    E = \frac{1}{2}m \left(r \dot{\theta} \hat{\underline{\theta}} + \dot{r} \hat{\underline{r}}\right)^2 + V(r)
\]
Noting that a unit vector dotted with themselves is one, and dotted with the (perpendicular) other gives zero:
\[
    E = \frac{1}{2}m \left(r^2 \dot{\theta}^2 + \dot{r}^2\right) + V(r)
\]
Where all the cross terms become zero. We can eliminate $\dot{\theta}$ by using:
\[
    l = mr^2 \dot{\theta} \implies \dot{\theta} = \frac{l}{mr^2}
\]
Hence:
\[
    E = \frac{1}{2}m \left(r^2 \frac{l^2}{m^2 r^4} + \dot{r}^2\right) + V(r)
\]
\[
    \boxed{E = \frac{1}{2}\frac{l^2}{mr^2} + \frac{1}{2}m \dot{r}^2 + V(r)}
\]

Note that:
\begin{itemize}
    \item The first term is the kinetic energy of the tangential motion. In the case of a circle, this gives $0.5 l^2/I_0 = 0.5 I^2 \omega^2 / I = \frac{1}{2}I \omega^2$ which is as expected.
    \item The second term is zero for circular motion where the radius is constant.
\end{itemize}
